%Dokumentenklasse
\documentclass[
	a4paper,										%Format
	12pt,											%Schriftgröße
	twoside,										%Zweiseitig
	headsepline,								%horizontale Linie oben
	footsepline,									%horizontale Linie unten
	captions=tableheading,					%Überschrift über Tabelle (nicht unter)
]{scrartcl}

\renewcommand{\thesubsection}{\thesection.\arabic{subsection}.}
\renewcommand{\thesubsubsection}{\thesection.\arabic{subsection}.\arabic{subsubsection}. \hspace{2pt}}
\renewcommand{\theparagraph}{\thesection.\arabic{subsection}.\arabic{subsubsection}.\arabic{paragraph}.}

%Pakete:
\usepackage[exponent-product = \cdot, output-decimal-marker = {,}]{siunitx}
\usepackage{textalpha}
\usepackage{floatrow}
\floatsetup[table]{style=plaintop}
\floatsetup[longtable]{LTcapwidth=table}
%\usepackage{overcite}
\usepackage{siunitx}
\usepackage[ngerman]{babel} 			%Neue deutsche Rechtschreibung
\usepackage{grffile}
\usepackage{graphicx} 					%Bilder einbinden
\usepackage{xcolor}						%Farben für Schrift und Plots
\usepackage{pgfplots}
\pgfplotsset{compat=1.18}					%Plots und Tabellen aus Datensätzen.
\usepackage{amsmath}						%Mathe
\usepackage{amssymb}					%Mathesymbole
\usepackage{longtable}					%Seitenübergreifende Tabellen zulassen
\usepackage{float}
\usepackage{tabularx}
\usepackage{geometry}					%Seitenränder
\usepackage{ae}								%skalierbare Schriften (im PDF-File)
\usepackage[activate]{microtype}		%optischer Randausgleich
\usepackage{setspace}					%Zeilenabstand: \onehalfspacing
\usepackage[utf8]{inputenc}				%Zeichenkodierung für Umlaute
\usepackage{scrlayer-scrpage}					%für Dokumentformatierung mit Pagestyle
\usepackage[version=4]{mhchem}						%Chemische Formeln (mit \ce )
\usepackage{chemfig}
\usepackage{siunitx}
\usepackage{caption}										%Überschriften formatieren
\captionsetup{labelfont={bf,it},tablename=Tab.,figurename=Fig.}				%Überschriften formatieren
\captionsetup[table]{singlelinecheck=false,justification=raggedright}
\captionsetup[figure]{justification=raggedright}				

%Abbildungs und Tabellennummerierung anpassen
\renewcommand*\figureformat {\figurename~\arabic{section}.\arabic{subsection}.\thefigure.}	%Punkt nach Abb. #
\renewcommand*\tableformat {\tablename~\arabic{section}.\arabic{subsection}.\thetable:}		%Doppelpunkt nach Tab. #
\renewcommand*\captionformat {~}								%Nur Abstand, kein Zeichen nach Tab. #: bzw Abb. #. 

%Literatur als Subsubsection definieren und von Literaturverzeichnis in Literatur umbenennen
\makeatletter
\renewcommand*\bib@heading{\subsubsection{\bibname}			
\@mkboth{\subsubsectionmarkformat \bibname}{\subsubsectionmarkformat \bibname}}
\makeatother
\addto\captionsngerman{\renewcommand*\bibname{Literature}} 

\usepackage{prettyref}						%Nummerierung der Tabellen und Abbildungen im Text an Tab. x.x.# anpassen
%%% Für Abbildungen %%%
\newrefformat{fig}{Fig.~\arabic{section}.\arabic{subsection}.\ref{#1}}
%%% Für Tabellen %%%
\newrefformat{tab}{Tab.~\arabic{section}.\arabic{subsection}.\ref{#1}}

%Links im Dokument:
\usepackage[									%als letztes Paktet einbinden!
	linkcolor=black,							%Farbe der Links
	citecolor=black,							%Farbe der Link-Ziele
	urlcolor=black,								%Farbe von URLs
	colorlinks=true,							%Links nicht als Kasten, sondern direkt als Text anzeigen
]{hyperref}										%Links im PDF-File, unterstützung von URLs mit /url{}

%Dokumentformatierungen
\setkomafont{sectioning}{\bfseries} 	%Serifenschrift in Überschriften
\setkomafont{subsection}{\LARGE}
\setkomafont{subsubsection}{\large}
\pagenumbering{arabic}					%arabische Seitenzahlen
\geometry{left=3cm, right=2cm, top=2.5cm, bottom=2cm} 	%Seitenränder
\parindent0pt 									%Absätze einrücken um 0Pt
\parskip 10pt									%Abstand nach Absatz
\onehalfspacing								%Zeilenabstand
\pagestyle{scrheadings}					%Kopf- und Fussleistenformatierung
\headsep0.4cm								%Abstand Kopfleiste-Text
\footskip0.8cm									%Abstand Fussleiste-Text
\urlstyle{same}								%URLs normal formatieren, nicht als Typewriter

%Befehle definieren

%%%%% create chemequation environment

\newcounter{chemequation}
\renewcommand*\thechemequation{(\Roman{chemequation})}
\newenvironment{chemequation}{%
	\stepcounter{chemequation}%
	\begin{equation}}%
{\tag*{\thechemequation}%
	\end{equation}}

%%%%% create chemtag to iterate and ref chemequations in an align environment

\newcommand*\chemtag{%
	\stepcounter{chemequation}%
	\tag*{\thechemequation}}

%%%%% to create a chemicalequation float

\usepackage{float}
\newfloat{chemicalequation}{htbp}{cex}%[chapter]
\floatname{chemicalequation}{}

\newcommand{\zitat}[1]{\textsuperscript{\cite{#1}}}

\newcommand{\bild}[4]{					%Bilder mit Beschriftung einfügen
	\begin{figure}[#4]						%#4: Platzierung (htbp)
		\centering								%Bilder mittig einfügen
		\captionbox								%Beschriftungsumgebung
			{#2\label{fig: #1}}				%#2: Beschriftung
			{\includegraphics[#3]{#1}}	%#1: Dateiname, #3: Bildeigenschaften
	\end{figure}} 

\newcommand{\formel}[2]{				%Formel einfügen und Referenz erzeugen
	\begin{equation}
		\label{eq: #1}							%#1: Formelname
		#2											%Formelinhalt
	\end{equation}}

\newcommand{\cformel}[2][]{				%Chemische Formel einügen und Referenz erzeugen
	\begin{chemequation}
		\label{ceq: #1}					%#1: Formelname (optional)
		\ce{#2}							%#2: Formelinhalt
	\end{chemequation}\\}

	% 2. Ein neuer, eigener Befehl NUR für Formeln mit Zusatz auf der rechten Seite:
\newcommand{\cformelpot}[3]{%
    \begin{chemequation}
        \label{ceq: #1}
        \ce{#2} \hspace{2cm} #3
    \end{chemequation}%
}

\newcommand{\ox}[2]{
	\overset{\text{#2}}{\ce{#1}}}
	
\newcommand{\tabelle}[6]{				%Tabelle mit Beschriftung einfügen
\begin{table}[#3]							%#3: Platzierung (htbp)
\centering										%Tabelle mittig einfügen
\captionbox										%Beschriftungsumgebung
{#2												%#2: Tabellenbeschriftung
\label{tab: #1}}								%#1: Tabellenname
{\begin{tabular}{#4}						%#4: Anzahl und Ausrichtung Spalten (l c r)
\hline 											%horizontale Linie
#5													%#5: Obere Zeile
\hline
#6													%#6: Datenzeilen
\hline 
\end{tabular}}
\end{table}}

\usepackage{array,ragged2e}
\newcolumntype{P}[1]{>{\RaggedRight\arraybackslash}p{#1}} %Verhindert ungewollte Lücken in Tabellen

\newcommand{\langetabelle}[6]{  %Seitenübergreifende Tabelle, Kopfzeile wird 
	\begin{centering}           %auf jeder neuen Seite ebenfalls angezeigt
		{\begin{longtable}[#3]{#4}
				\caption{#2\label{tab: #1}}\\				
				\hline
				#5
				\hline
				\endfirsthead
                \caption{#2\label{tab: #1}}\\
				\hline
				#5
				\hline
				\endhead
				#6												
				\hline 
		\end{longtable}}
\end{centering}}

%Plots erstellen&&&&&&&&

%#1: Referenz
%#2: Name
%#3: x label
%#4: y label
%#5: x-werte,y-werte (Spaltenweise)
%#6: Quelle (.txt Datei)
%#7: "Nummerierung" (manuell) <.<

\newcommand{\plot}[6]{
	\begin{figure}[H]	
		\centering
		\begin{tikzpicture}
	
		\begin{axis}[
		title style={at={(0.5,0)},
		anchor=north,
		yshift=-40},
		xlabel={#3},
		ylabel={#4},
		legend pos=south east,
		legend entries={},
		]
		\addplot[color=black,mark=*,only marks] table [#5] {#6};
		\label{plot: plot1}
		\end{axis}
		\end{tikzpicture}
		\caption{#2}
		\label{fig: #1}
	\end{figure}
}

\newcommand{\molL}[1]{\ensuremath{\mathrm{\SI{#1}{\frac{mol}{L}}}}} % \molL{1} => \SI{1}{\frac{mol}{L}}
\newcommand{\mL}[1]{\ensuremath{\mathrm{\SI{#1}{mL}}}} % \mL{1} => \SI{1}{mL}

\newcommand{\Rn}[1]{\romannumeral#1\relax}                   %\Rn{1} = i
\newcommand{\RN}[1]{\uppercase\expandafter{\romannumeral#1}} %\RN{1} = I

%Dokument:%%%%%%%%%%%%%%%%%%%%%%%%%%%%%%%%%%%%%%%%%%%%%%%%%%%%%%%%%%%%%%%%%%%%%%%%%%%%
\begin{document}
	
	%Titelseite
	\begin{titlepage}
		\begin{center}
			\hrule													%waagerechte Linie
			\vspace*{5cm}
			\Huge {\textbf {Experimental Series 7}}\\						%Titel Versuchsreihe
			\vspace*{0.5cm}
			%		\large {\bf {Untertitel}}\\						%Untertitel
			\vspace*{1cm}
			\normalsize
			Lab Report	\\									
			Nanoscience Labartory Course \\ for Nanoscience Students\\
			\vspace*{0.5cm}
			\Large {\textbf{Maxim Popow, Michael Quint}}		\\						%Name
			\vspace*{0.5cm}
					\normalsize{Seats 8b, 9a}								%Platznummer
			\vfill
		\end{center}
		\begin{tabbing}
			\hspace{3cm}\=\kill
			Assistent: \> Helena Casara \\						% Name Assistent
			Version: \> 1 \\									%Version
			Date: \> 15.11.2026 \\								%Datum
		\end{tabbing}
		\hrule														%Waagerechte Linie
\end{titlepage}

\setcounter{page}{1}
\setcounter{section}{7}								%Versuchstag #.x.x bzw Versuchsreihe
\setcounter{subsection}{0}							%Versuchsnummer x.#.x startet bei 0 für 1 
\setcounter{secnumdepth}{4}
	
\subsection{Conductivity of Solutions}

\subsubsection{Introduction and Objectives}

The specific conductivity of several solutions was to be determined.
The molar conductivity as well as the equivalent conductivity of sodium chloride $(1\ \mathrm{\frac{mol}{L}})$ and magnesium sulfate $(1\ \mathrm{\frac{mol}{L}})$ solutions
were to be calculated.

\subsubsection{Theoretical Background}

Solutions are defined as homogenous mixtures.
The kind of mixtures, that will be discussed during the following text, are in a liquid state of matter,
since they are mostly in that form and thus are the most important.
Whenever there is an excess of a substance in a solution, that substance is defined as a solvent.
Secondary substances are called solutes.

\ce{H2O} is a typical solvent. It is simultaneously a polar solvent, due to its partial charge between the oxygen atom and the hydrogen atoms.
Many chemical compounts in water can dissolve (or react) into differently charged ions, such as most salts (\ce{MX}), acids (\ce{HA}) and bases (\ce{B}).
\cformel[a1]{MX_{(s)} <=>[H2O][dissociation] M^{+}_{(aq)} + X^{-}_{(aq)}}
\cformel[a2]{HA_{(aq)} + H2O_{(l)} <=>[][protolysis] A^{-}_{(aq)} + H3O^{+}_{(aq)}}
\cformel[a3]{B_{(aq)} + H2O_{(l)} <=> HB^{+}_{(aq)} + OH^{-}_{(aq)}}

These solutes are being hydrated,
which means that the charged ions form a hydration shell with the opposite charged side of the solvent.\zitat{kopie1}

\bild
		{NaCl.png}					%Dateiname des Bildes, eindeutiger Name des Bildes, Verweis: \prettyref{fig: Bildname}
		{Dissociation of sodium chloride into hydrated ions.%Text, der unter dem Bild erscheint
		\zitat{kopie2}}			%Name der Quelle im Literaturverzeichnis
		{scale=0.5}		%angle, draft, scale, height, width
		{h}				%h: here, t: top, b: bottom, p: page; bestimmt, wo das Bild eingefügt wird.

If hydrated ions are involved in a solution, that solution is called an electrolyte.
Within these electrolytes hydrated ions are able to move more freely
and since they carry a non-zero charge,
an electric current $I$ is flowing,
whenever the rate of change of the electric charge $Q$ is not zero
i.e. when a positve anode and a negative cathode are placed in the solution.
\formel{streber1}{I = \frac{dQ}{dt}}

When there is a flowing current, Ohm's law applies, where $U$ is defined as the voltage and $R$ as the resistance:
\formel{streber2}{U = R \cdot I}
$R$ is dependent on the specific resistance $\rho$, while $A$ is the area of the two electrodes and $d$ the distance between the electrodes.
\formel{streber3}{R = \rho \cdot \frac{d}{A}}
The conductance $L$ and specific conductivity $\kappa$ are defined as the reciprocals of $R$ or rather $\rho$\zitat{kopie3}:
\formel{streber4}{L = \frac{1}{R}}
\formel{streber5}{\kappa = \frac{1}{\rho}}

The specific conductivity can now be used to introduce the molar conductivity $\Lambda_{\mathrm{m}}$:
\formel{streber6}{\Lambda_{\mathrm{m}} = \frac{\kappa}{c}}
Since some ions hold a different charge number it is sensible to also introduce the equivalent conductivity $\Lambda_{\mathrm{eq}}$,
where $z$ is defined as the charge number of the ion:
\formel{streber7}{\Lambda_{\mathrm{eq}} = \frac{\Lambda_{\mathrm{m}}}{z} = \frac{\kappa}{c \cdot z}}

As Ostwald's dilution law states, the degree of dissociation $\alpha$ of weak electrolytes is proportional to the square root of the reciprocal of the concentration of the solute.
\formel{streber8}{\alpha \propto \frac{1}{\sqrt{c}}}

A higher concentration means however, that the solutes interact with each other, lowering the 'effective' concentration.
In these cases, the activity $a$ is introduced, which itself is proportional to its activity coefficient $\gamma_{\pm}$.
The activity coefficient effectivly measures the degree of 'ideality' by assuming a value between 0 and 1, whereas a $\gamma_{\pm}$ of 1 means 'fully ideal'.
$c^{\circ}$ is defined as the standard concentration of $1\ \mathrm{\frac{mol}{L}}$:
\formel{a1}{a = \gamma_{\pm} \cdot \frac{c}{c^{\circ}}}
An ideal solution can be achieved by fully diluting it.
The molar limiting conductivity $\Lambda^{\circ}_{\infty}$ is therefore introduced as the molar conductivity of an infinitly diluted solution.
As per eq. \prettyref{eq: streber8}, the molar conductivity for weak electrolytes can be described as\zitat{kopie4}:
\formel{a2}{\Lambda_{\mathrm{m}} = \alpha \cdot \Lambda^{\circ}_{\infty}}
Analogous to weak electrolytes, Kohlrausch's square root law applies to strong electrolytes for $c < 0.1 \ \mathrm{\frac{mol}{L}}$, whereas $K$ is the Kohlrausch constant\zitat{kopie5}:
\formel{a3}{\Lambda_{\mathrm{m}} =\Lambda^{\circ}_{\infty} - K \cdot \sqrt{c}}
It is valuable to mention that solutions involving hydronium (\ce{H3O+}) and hydroxide (\ce{OH-}) ions possess a high specific conductivity,
since the charge transport does not include the whole ion, but rather the proton (\ce{H+}) commits a 'change of place' towards other molecules,
resulting in a relatively low resistance.\zitat{kopie4}

\subsubsection{Experimental Section}

The specific conductivities of several solutions with a concentration of 0.01 $\mathrm{\frac{mol}{L}}$ were measured with a combination measuring instrument,
as represented in \prettyref{tab: LF1}.

\langetabelle
	{LF1}
	{Measured specific conductivities of solutions.}
	{H}
	{p{5cm} p{5cm}}
	{Solution (0.01 $\mathrm{\frac{mol}{L}}$) & $\kappa$ / ${\frac{\mu\mathrm{S}}{\mathrm{cm}}}$ \\}
	{
	Glucose
	& 20.50 \\

	Urea
	& 22.60 \\

	\ce{CH3COOH}
	& 177.0 \\

	\ce{NH3}
	& 79.20 \\

	\ce{NaCH3COO}
	& 955.0 \\

	\ce{NaCl}
	& 985.0 \\

	\ce{HCl}
	& 3850 \\

	\ce{NaOH}
	& 2150 \\

	\ce{MgSO4}
	& 2440 \\

	
	}

Later on, the specific conductivities of \ce{NaCl, CH3COOH and NaCH3COO} solutions were measured with a combination measuring instrument,
as represented in \prettyref{tab: LF2}; however this time with a concentration of 0.05 $\mathrm{\frac{mol}{L}}$.

\langetabelle
	{LF2}
	{Measured specific conductivities of solutions.}
	{H}
	{p{5cm} p{5cm}}
	{Solution (0.05 $\mathrm{\frac{mol}{L}}$) & $\kappa$ / ${\frac{\mu\mathrm{S}}{\mathrm{cm}}}$ \\}
	{ 
	
	\ce{CH3COOH}
	& 388.0 \\

	\ce{NaCH3COO}
	& 5280 \\

	\ce{NaCl}
	& 5580 \\
	
	}

Lastly, the specific conductivity of ammonia (\ce{NH3}) and acetic acid (\ce{CH3COOH}) solutions were measured, along with a 1:1 mixture of said solutions.

\langetabelle
	{LF3}
	{Measured specific conductivities of solutions.}
	{H}
	{p{5cm} p{5cm}}
	{Solution (0.05 $\mathrm{\frac{mol}{L}}$) & $\kappa$ / ${\frac{\mu\mathrm{S}}{\mathrm{cm}}}$  \\}
	{
	
	\ce{NH3}
	& 380.0 \\

	\ce{CH3COOH}
	& 388.0 \\

	\ce{NH3 + CH3COOH}
	& 1884 \\
	
	}

\subsubsection{Results and Discussion} 

The measured specific conductivities are roughly compatible with the values that could have been predicted.
Glucose and urea do not dissolve easily in demineralised water,
thus their specific conductivity is approximately on par with pure demineralised water.

Ammonia (\ce{NH3}) and acetic acid (\ce{CH3COOH}) are weak bases/acids, which results in next to zero deprotolysis/dissociation and a small conductivity.

Hydrochloric acid (\ce{HCl}) and sodium hydroxide (\ce{NaOH}) however nearly fully dissolve in water.
The dissolution also results in hydronium and hydroxide ions being involved, further increasing the specific conductivities.

Sodium acetate (\ce{NaCH3COO}), sodium chloride (\ce{NaCl}) and magnesium sulfate (\ce{MgSO4}) also fully dissolve in water,
thus are strong electrolytes, but since each magnesium and sulfate ion carries twice the charge of a sodium or a chloride ion,
their measured conductivity should also be substantually larger.\\
The molar conductivity of sodium chloride and magnesium sulfate can now be calculated as shown in eq. \prettyref{eq: streber6}:

\formel{b1}{\Lambda_{\mathrm{m, \ce{NaCl}}} = \frac{\kappa_{\ce{NaCl}}}{c} = \frac{985.0\ \mathrm{\frac{\mu S}{cm}}}{0.01\ \mathrm{\frac{mol}{L}}} = 98.5\ \mathrm{\frac{\mu S \cdot cm^2}{mol}}}
\formel{b2}{\Lambda_{\mathrm{m, \ce{MgSO4}}} = \frac{\kappa_{\ce{MgSO4}}}{c} = \frac{2440\ \mathrm{\frac{\mu S}{cm}}}{0.01\ \mathrm{\frac{mol}{L}}} = 244\ \mathrm{\frac{\mu S \cdot cm^2}{mol}}}

The equivalent conductivity as per eq. \prettyref{eq: streber7} is as follows:

\formel{b3}{\Lambda_{\mathrm{eq, \ce{NaCl}}} = \frac{\Lambda_{\mathrm{m, \ce{NaCl}}}}{z} = \frac{98.5\ \mathrm{\frac{\mu S \cdot cm^2}{mol}}}{1} = 98.5\ \mathrm{\frac{\mu S \cdot cm^2}{mol}}}
\formel{b4}{\Lambda_{\mathrm{eq, \ce{MgSO4}}} = \frac{\Lambda_{\mathrm{m, \ce{MgSO4}}}}{z} = \frac{244\ \mathrm{\frac{\mu S \cdot cm^2}{mol}}}{2} = 122\ \mathrm{\frac{\mu S \cdot cm^2}{mol}}}

If the concentration of a weak acid (such as acetic acid) is quintupled, its degree of dissociation $\alpha$ gets smaller, resulting in 'only' double the specific conductivity,
as seen in \prettyref{tab: LF2}.
But since sodium chloride and sodium acetate fully dissolve in water, their specific conductivity also roughly quintuples.

\formel{c1}{ \frac{\kappa_{\ce{NaCl}, 2}}{\kappa_{\ce{NaCl}, 1}} = \frac{5580\ \mathrm{\frac{\mu S}{cm}}}{985.0\ \mathrm{\frac{\mu S}{cm}}} = 5.665 > 5}
\formel{c2}{ \frac{\kappa_{\ce{NaCH3COO}, 2}}{\kappa_{\ce{NaCH3COO}, 1}} = \frac{5280\ \mathrm{\frac{\mu S}{cm}}}{955.0\ \mathrm{\frac{\mu S}{cm}}} = 5.529 > 5}

But as per Kohlrausch's square root law in eq. \prettyref{eq: a3} the ratio of the two specific conductivities should be below 5,
because the molar conductivity gets lower, the higher the concentration is.
This discrepancy can be attributed to measurement uncertainties and unclean containers, where the solutions were stored.

Ammonia and acetic acid by themselves are weak electrolytes, since they do not fully undergo a complete (de-)protolysis.
If mixed together, an acid-base reaction takes place:
\cformel[d1]{NH3 + CH3COOH <=> NH4+ + CH3COO-}
In contrast, ammonium acetate (\ce{NH4CH3COO}) fully dissolves in water, resulting in a relatively high specific conductivity.

\subsubsection{Summary}
The measured specific conductivities of the solutions were verified for consistency.
The effect of strong and weak electrolytes can be explained through Ostwald's dilution law,
with relatively moderate discrepancies regarding Kohlrausch's square root law.
The specific conductivity of the mixture of two weak electrolytes can also be explained by an acid-base reaction.
The molar and equivalent conductivities of sodium chloride and magnesium sulfate were also successfully calculated.

\subsubsection{Safety and Disposal}
		
The substances used in this experiment with danger symbols, H- and P-statements and references regarding their disposal are cited in \prettyref{tab: Chemikalien}.\zitat{test} Used CMR-substances are listed in \prettyref{tab: KMR_leer}.\zitat{test}
	
\langetabelle
	{Chemikalien}
	{Used chemicals.}
	{H}  % hier ggf. ein großes H eintragen, wenn die Tabelle sonst im Anhang verschwindet
	{l p{1.9cm} p{2.5cm} p{5cm} c}
	{Substance & Symbol & H-Statements & P-Statements & Disposal \\}
	{ % Erster Stoff

	\ce{NaOH\ (1\ \mathrm{\frac{mol}{L}})}
		& \begin{tabular}[t]{l} GHS05 \end{tabular}
		& 290, 314
		& 260, 264, 280, \newline 301+330+331, \newline 303+361+353, \newline 304+340, \newline 305+351+338, 310, \newline 321, 405, 501
		& (1) \\

	\ce{CH3COOH}
		& \begin{tabular}[t]{l} GHS02 \\ GHS05 \end{tabular}
		& 226, 290 \newline 314
		& 210, 233, 240, 241, \newline 242, 243, 260, 264, \newline 280, 301+330+331, \newline 303+361+353, \newline 304+340, \newline 305+351+338, 310, \newline 321, 370+378, \newline 403+235, 405, 501
		& (2) \\

	\ce{HCl\ (1\ \mathrm{\frac{mol}{L}})}
		& \begin{tabular}[t]{l} GHS02 \\ GHS05 \end{tabular}
		& 290, 314 \newline 335
		& 260, 280, 390, \newline 303+361+353, \newline 305+351+338, 310, \newline 312
		& (2) \\

	\ce{NH3\ (25\ \%)}
		& \begin{tabular}[t]{l} GHS02 \\ GHS05 \end{tabular}
		& 290, 314 \newline 335, 400
		& 260, 264, 273, 280, \newline 301+330+331, \newline 303+361+353, \newline 304+340, \newline 305+351+338, 310, \newline 321, 391, 405, 501
		& (3) \\

	}

\langetabelle
	{KMR_leer}
	{Carcinogenic, mutagenic or reprotoxic hazardous substances used (CMR).}
	{H}
	{l P{3.8cm} P{2.4cm} P{3.5cm} c}
	{Substances & H-Statements & CMR-Classification & Used Quantity & Disposal \\}
	{ - & - & - & - & - \\}
	
Disposal: \\
\begin{itemize}
	\item[(1)] The solution was disposed of in the container for other bases. 
	\item[(2)] The solution was disposed of in the container for other acids.
	\item[(3)] The solution was acidified and then disposed of in the container for other acids.
\end{itemize}	

\begin{thebibliography}{4} %manuelles Quellenverzeichnis

	\bibitem{kopie1}
						{E. Riedel, C. Janiak}, 
			\textit	{Anorganische Chemie}, 
						{10. Auflage}, 
						{de Gruyter}, 
						{Berlin}, 
			\textbf	{2022}, 
			S.			{336}.

	\bibitem{kopie2}
						{E. Riedel, C. Janiak}, 
			\textit	{Anorganische Chemie}, 
						{10. Auflage}, 
						{de Gruyter}, 
						{Berlin}, 
			\textbf	{2022}, 
			S.			{337}.

	\bibitem{kopie3}
						{E. Riedel, C. Janiak}, 
			\textit	{Anorganische Chemie}, 
						{10. Auflage}, 
						{de Gruyter}, 
						{Berlin}, 
			\textbf	{2022}, 
			S.			{338}.
	
	\bibitem{kopie4}
						{E. Riedel, C. Janiak}, 
			\textit	{Anorganische Chemie}, 
						{10. Auflage}, 
						{de Gruyter}, 
						{Berlin}, 
			\textbf	{2022}, 
			S.			{339-341}.

	\bibitem{kopie5}
						{P. Kurzweil}, 
			\textit	{Chemie: Grundlagen, technische Anwendungen, Rohstoffe, Analytik und Expermiente}, 
						{11. Auflage}, 
						{Springer Vieweg}, 
						{Wiesbaden}, 
			\textbf	{2020}, 
			S.			{222}.


	\bibitem{test}
						{DGUV}, 
						„{GESTIS-Stoffdatenbank}“, found at
			\url		{https://gestis.dguv.de/}, 
			\textbf	{01.11.2026}.
			
\end{thebibliography}

%---------------------------------------------------------------------------------------------------------------------------------
\newpage
	
\subsection{V 7.2.}

\subsubsection{Introduction and Objectives}

\subsubsection{Theoretical Background}

\subsubsection{Experimental Section}

\subsubsection{Results and Discussion} 

\subsubsection{Summary}

\subsubsection{Safety and Disposal}
		
The substances used in this experiment with danger symbols, H- and P-statements and references regarding their disposal are cited in \prettyref{tab: Chemikalien3}.\zitat{hauke3} Used CMR-substances are listed in \prettyref{tab: KMR_leer3}.\zitat{hauke3}
	
\langetabelle
	{Chemikalien3}
	{Used chemicals.}
	{H}  % hier ggf. ein großes H eintragen, wenn die Tabelle sonst im Anhang verschwindet
	{l p{1.9cm} p{2.5cm} p{5cm} c}
	{Substance & Symbol & H-Statements & P-Statements & Disposal \\}
	{ % Erster Stoff

	\ce{NaOH}
		& \begin{tabular}[t]{l} GHS05 \end{tabular}
		& 290, 314
		& 234, 260, 280, \newline 303+361+353, \newline 304+340+310, \newline 305+351+338,
		& (1) \\

    }

\langetabelle
	{KMR_leer3}
	{Carcinogenic, mutagenic or reprotoxic hazardous substances used (CMR).}
	{H}
	{l P{3.8cm} P{2.4cm} P{3.5cm} c}
	{Substances & H-Statements & CMR-Classification & Used Quantity & Disposal \\}
	{ - & - & - & - & - \\}
	
Disposal: \\
\begin{itemize}
    \item[(1)] 
\end{itemize}	


\begin{thebibliography}{4}		%manuelles Quellenverzeichnis


	\bibitem{hauke3}
						{DGUV}, 
						„{GESTIS-Stoffdatenbank}“, found at 
			\url		{https://gestis.dguv.de/}, 
			\textbf	{01.11.2026}.
	

\end{thebibliography}

\end{document}

%\newpage

%%%%%%%%%%%%%%%%%%%%%%%%%%%%%%%%%%%%%%%%%%%%%%%%%%%%%%%%%%%%%%%
Neuer Versuch:
\subsection{Versuchsname}

\subsubsection{Einleitung und Aufgabenstellung}
\subsubsection{Theoretischer Hintergrund}
\subsubsection{Experimenteller Teil}
\subsubsection{Ergebnisse und Diskussion}
\subsubsection{Zusammenfassung}
\subsubsection{Sicherheit und Entsorgung}	
Die für den Versuch verwendeten Chemikalien sind in \prettyref{tab: entsorgung1} mit Gefahrensymbolen, H- und P-Sätzen sowie Hinweisen zur Entsorgung aufgeführt.\zitat{Gefahren1} Verwendete KMR-Stoffe sind in \prettyref{tab: KMR1} aufgeführt.
	
\langetabelle
	{entsorgung1}
	{Verwendete Chemikalien.}
	{H}  % hier ggf. ein großes H eintragen, wenn die Tabelle sonst im Anhang verschwindet
	{l p{1.4cm} p{2.2cm} p{4.5cm} c}
	{Substanz & Symbol & H-Sätze & P-Sätze & Entsorgung \\}
	{\\}
    

\langetabelle
	{KMR1}
	{Verwendete kanzerogene, mutagene oder reproduktionstoxische Gefahrstoffe (KMR).}
	{H}
	{l P{3.8cm} P{2.4cm} P{3.5cm} c}
	{Substanz & H-Sätze & KMR-Einstufung & Verwendete Menge & Entsorgung \\}
	{ - & - & - & - & - \\}
Entsorgung: \\
\begin{itemize}
    \item[(1) Die verwendete Lösung wurde mit Salzsäure angesäuert und im Behälter für nicht-nitrathaltige Säuren entsorgt.] 
\end{itemize}	

\newpage

\begin{thebibliography}{4}		%manuelles Quellenverzeichnis
	\bibitem{Gefahren1}
						{H. Heller}, 
			\textit	{Grundlagenpraktikum der Chemie für Studierende der Nanowissenschaften}, 
						{Version 1.18},
                        {Universität Hamburg},
						{Hamburg}, 
			\textbf	{2023}, 
			S.{34}.
			
			
\end{thebibliography} 
\newpage

Nützliche Beispiele:

%Tabelle für Reagenzgläserinhalte
\tabelle
	{Tabellenname}												%Eindeutige Bezeichnung der Tabelle. \prettyref{tab: Tabellenname}
		{Tabellenüberschrift.}	%Inhalt der Reagenzgläser ... bis...
		{htb}	
		{l l l}					
		{Reagenzglas & Inhalt & Beschreibung}				
		{1 & 2~mL NaOH & klare, farblose Lsg.\\
		2 & & 
		}										


%Formel Löslichkeit KClO4 -> chemische Formel
\begin{align}
    \chemtag
    \ce{KClO4(aq) <=> K+(aq) + ClO4-(aq)}
\end{align}

%hier allgemeine Formel 
\formel
		{formel2}				%eindeutige Bezeichnung der Formel, Verweis: \ref{eq: formelname}
		{L_{\mathrm{A}_{a} B_{b}} = \sqrt[m+n]{\frac{K_{\mathrm{L}}}{a^{a} \cdot b^{b}}}  }


Bilder:
	\bild
		{bildname}					%Dateiname des Bildes, eindeutiger Name des Bildes, Verweis: \prettyref{fig: Bildname}
		{Bildbeschreibung.		%Text, der unter dem Bild erscheint
		\zitat{Quelle}}			%Name der Quelle im Literaturverzeichnis
		{Bildeigenschaften}		%angle, draft, scale, height, width
		{Platzierung}				%h: here, t: top, b: bottom, p: page; bestimmt, wo das Bild eingefügt wird. p: spezielle Seite für Abbildungen und Tabellen, Kombinationen möglich, Versuch der Platzierung gemäß Reihenfolge.
	
		
	Bildeigenschaften:
		angle - Bestimmt den Drehwinkel um einen Referenzpunkt für die Grafik. Positive Werte drehen die Grafik gegen den Uhrzeigersinn, negative mit ihm. 
		draft - Ist draft gesetzt, wird nur ein Platzhalter der Grafik angezeigt. Beschleunigt den Übersetzungsvorgang des Dokuments erheblich. 
		scale - Skaliert die Grafik anhand eines Skalierungsfaktors. 
		height - Skaliert die Grafik auf die angegebene Höhe. Die Angabe der Größe muss in einer LaTeX-spezifischen Einheit erfolgen. 
		width - Skaliert die Grafik auf die angegebene Breite. Die Angabe der Größe muss in einer LaTeX-spezifischen Einheit erfolgen. 	
	

Formeln:
	\formel
		{formelname}				%eindeutige Bezeichnung der Formel, Verweis: \ref{eq: formelname}
		{Formel}						%


Tabellen:
	\tabelle
		{Tabellenname}												%Eindeutige Bezeichnung der Tabelle. \prettyref{tab: Tabellenname}
		{Tabellenüberschrift.\zitat{Quelle}}						%Text, der über der Tabelle erscheint, Quelle der Daten
		{Platzierung}													%h: here, t: top, b: bottom, p: page; bestimmt, wo das Bild eingefügt wird. p: spezielle Seite für Abbildungen und Tabellen, Kombinationen möglich, Versuch der Platzierung gemäß Reihenfolge.		
		{Spaltenanzahl_und_Ausrichtung}						%{ l c r }: 3 Spalten, links mittig und rechts; { c c c c }: 4 Spalten, alle mittig
		{Titelzeile}														%Obere Zeile(n), danach folgt ein waagerechter Strich, Spalten durch & trennen, Neue Zeile durch \\
		{Datenzeilen}													%Zeilen unter dem waagerechten Strich, Spalten durch & trennen, Neue Zeile durch \\
		
	Zellen verbinden:	
		\multicolumn{Anzahl_der_verbunden_Spalten}{Ausrichtung (l/c/r}{Inhalt}
	
Reaktionen:
Formeln: 
\begin{center}
\ce{NH+4}, \ce{Fe3(CO)12}, \ce{Al^{3+}},\ce{Cr2O^{2-}7}
\end{center}

Isotope:
\begin{center}
\ce{^{37}_{17}Cl}, \ce{^{32}_{16}S^{2-}}
\end{center}

Bindungen: 
\begin{center}
\ce{H-O-N=S}, \ce{Cl-C#N}
\end{center}

Reaktionsgleichungen: 
\begin{center}
\ce{2H2 + O2 -> 2H2O}
\ce{MgCl2 ->[\text{Elektrolyse}] Mg\,(l) + Cl_{2}\,(g)}
\end{center}

Redoxreaktion:
\begin{center}
$\overset{0}{\ce{S}} + \overset{0}{\ce{O2}} \ce{->} \overset{+\text{IV}}{\ce{S}}+\overset{-\text{II}}{\ce{O2}}$
\end{center}


Literaturverweise:
		\bibitem{zeitschrift}
						{A. Autor}, 
			\textit	{Zeitschrift} 
			\textbf	{Jahr}, 
			\textit	{Band} 
						{Seite-Seite}.
		
		\bibitem{buch_ohne_herausgeber}
						{A. Autor}, 
			\textit	{Titel}, 
						{Verlag}, 
						{Ort}, 
			\textbf	{Jahr}, 
			S.			{Seite-Seite}.
			
		\bibitem{buch_mit_herausgeber}
						{A. Autor}, 
			\textit	{Titel}, 
						{Auflage}, 
						{Verlag}, 
						{Ort}, 
			\textbf	{Jahr}, 
			S.			{Seite-Seite}.
			
		\bibitem{dissertation}
						{A. Autor}, 
						{Dissertation}, 
						{Universität}, 
			\textbf	{Jahr}.
			
		\bibitem{internetseite}
						{A. Autor}, 
						„{Titel}“, zu finden unter 
			\url		{www.xyz.de}, 
			\textbf	{Jahr}.

% Kommentare Sophie:

% <- Für Kommentare oder Code der gerade nicht gebraucht wird, wird beim kompilieren nicht beachtet
%Um Bilder einzufügen: \bild{'Dateiname' (Datei muss im gleichem Ordner mit der TeX Datei liegen)}{'Bildunterschrift'}{width= ...cm (Größe des Bildes)}{H (Ort des Bildes, H=Hier)}

%Für Summenformeln im Text: \ce{z.B.: H2SO4 (Zahlen werden automtisch tief gestellt)}

% Um Formeln im Text zu referieren: \prettyref{fig (für Bilder und Graphen)/ eq (für mathematische Formeln)/ ceq (für chemische Formeln): 'Referenz'}

%Für chemische Formeln: \cformel['Referenz']{2\ox{Fe}{+\RN{3}}Cl3 + \ox{Fe}{+\RN{2}}Cl2 + 8NH4+ + 8 OH- \ce{<=>[{N2H4}]} \ox{Fe}{+\RN{2}}O\cdot \ox{Fe}{+\RN{3}}2O3 \ce{v} + 4 H2O + 8NH4+ + 4 Cl-}

%Für mathematische Formeln: \formel{'Referenz'}{2\cdot d\cdot \ce{sin}(\theta)= n \lambda}

%Tabellen: Online Converter um Exel Tabellen in LaTeX Format zu bringen: https://www.tablesgenerator.com
%\tabelle{'Referenz'}{'Tabellenüberschrift'}{H}{'Breite der Spalte bzw. Textausrichtung' (Spaltenbreite passt sich Textlänge an: l=links, r=rechts, c=zentriert. Spaltenbreite fest und Blocksatz: p{...cm}), Jede Spalte muss einzeln definiert werden, z.B l l l = 3 Spalten}{'Spaltenüberschrift' muss der Anzahl an Spalten entsprechen. Spaltenüberschriften müssen durch ein & getrennt werden. Außerdem wichtig am Ende: \\}{'Tabelleneinträge', müssen ebenfalls durch ein & getrennt werden und am Ende der Zeile ein \\ haben
%	- & - & - & \\
%	- & - & - & \\}
%Beispiel für eine Tabelle (zum Ausführen '%' entfernen):
%\tabelle{Auflösungsvermögen}{Experimentell ermitteltes und theoretisches Auflösungsvermögen für Objektiv 20X und 63X bei schräger und gerader Beleuchtung.}{H}{l l p{2cm} p{3cm}  p{3cm}}{Objektiv&Beleuchtung& Spalt mit einem Peak&experimentelles Auflösungsvermögen~/~nm&theoretisches Auflösungsvermögen~/~nm\\}
%	{20X&gerade&16&>750&1125\\
%	20X&schräg&9&690&563\\
%	63X&gerade&11&>750&600\\
%	63X&schräg&7&584&300\\}